% !TeX root = bbchallenge-paper.tex

\documentclass[sigconf,screen]{acmart}

%% Rights management information.
%% TODO: Update these values when you receive the rights form from ACM.
\setcopyright{acmlicensed}
\copyrightyear{2026}
\acmYear{2026}
\acmDOI{XXXXXXX.XXXXXXX}

%% STOC 2026 conference information
\acmConference[STOC '26]{Proceedings of the 58th Annual ACM Symposium on Theory of Computing}{June 22--26, 2026}{Salt Lake City, UT, USA}
\acmISBN{978-1-4503-XXXX-X/26/06}

\usepackage{floatpag}
\usepackage{subcaption}
\usepackage{tikz}

\usepackage{wrapfig}
\usepackage{refcount}
\usetikzlibrary{calc}
\usetikzlibrary{arrows.meta}

\usepackage{algorithm}
\usepackage[noend]{algpseudocode}
\usepackage{colortbl}

\definecolor{graySymb1}{RGB}{85,85,85}
\definecolor{graySymb2}{RGB}{170,170,170}

\usepackage{bold-extra}% bold texttt

\usepackage{mathtools}

\theoremstyle{definition} % don't use italics
\newtheorem{theorem}{Theorem}[section]
\newtheorem{conjecture}{Conjecture}[section]
\newtheorem{problem}{Problem}[section]
\newtheorem{definition}{Definition}[section]
\newtheorem{lemma}{Lemma}[section]
\newtheorem{proposition}{Proposition}[section]
\newtheorem{corollary}{Corollary}[section]
\numberwithin{equation}{section}


\theoremstyle{definition} % emphasize the "Remark N." with italics, not bold
\newtheorem{observation}{Observation}[section]
\newtheorem{example}{Example}[section]
\newtheorem{remark}{Remark}[section]

\usepackage{xassoccnt}
\DeclareCoupledCountersGroup{theorems}
\DeclareCoupledCounters[name=theorems]{theorem,definition,lemma,proposition,corollary,remark,observation,example}

\usepackage{xspace,wrapfig,multicol}
\usepackage[textsize=tiny,color=lightgray]{todonotes}
\usepackage[normalem]{ulem} % sout
\usepackage{stmaryrd}


\newcommand{\ts}[1]{{\color{red}#1}}
\newcommand{\tsi}[1]{\todo[inline]{TS: #1}}
\newcommand{\tsm}[1]{\todo{TS: #1}}
\newcommand{\tss}[2]{{\ts{\sout{#1}}} {\ts{#2}}}
\newcommand{\jb}[1]{{\color{blue}#1}}
\newcommand{\jbi}[1]{\todo[inline]{JB: #1}}
\newcommand{\jbm}[1]{\todo{JB: #1}}
\newcommand{\jbs}[2]{{\jb{\sout{#1}}} {\jb{#2}}}
\newcommand{\tabi}{\hspace{\algorithmicindent}}
\newcommand{\Lim}[1]{\raisebox{0.5ex}{\scalebox{0.8}{$\displaystyle \lim_{#1}\;$}}}
\newcommand{\N}{\mathbb{N}}
\newcommand{\Z}{\mathbb{Z}}

\newcommand{\lheadTM}[1]{\stackrel{#1}\triangleleft}
\newcommand{\rheadTM}[1]{\stackrel{#1}\triangleright}

\newcommand{\tm}[1]{\href{https://bbchallenge.org/#1}{\texttt{\nolinkurl{#1}}}}

\usepackage{adjustbox}


\definecolor{colorA}{RGB}{255,0,0}
\definecolor{colorB}{RGB}{255,128,0}
\definecolor{colorC}{RGB}{0,0,255}
\definecolor{colorD}{RGB}{0,255,0}
\definecolor{colorE}{RGB}{255,0,255}

\usepackage{epigraph}

\newcommand{\stateA}{{\textcolor{colorA}{A}}\xspace}
\newcommand{\stateB}{{\textcolor{colorB}{B}}\xspace}
\newcommand{\stateC}{{\textcolor{colorC}{C}}\xspace}
\newcommand{\stateD}{{\textcolor{colorD}{D}}\xspace}
\newcommand{\stateE}{{\textcolor{colorE}{E}}\xspace}

\newcommand{\stateAx}{{\textcolor{colorA}{A}}}
\newcommand{\stateBx}{{\textcolor{colorB}{B}}}
\newcommand{\stateCx}{{\textcolor{colorC}{C}}}
\newcommand{\stateDx}{{\textcolor{colorD}{D}}}
\newcommand{\stateEx}{{\textcolor{colorE}{E}}}

\usepackage{enumitem}

\newcommand{\szero}{\texttt{0}\xspace}
\newcommand{\sone}{\texttt{1}\xspace}


\newcommand{\BBtheSecondTNF}{61}
\newcommand{\BBtheThirdTNF}{5{,}417}

\newcommand{\BBtheFourth}{107}
\newcommand{\SigmaTheFourth}{13}
\newcommand{\BBtheFourthTNF}{858{,}909}

%used to be \hookrightarrow but savask says its to be mistaken with injection
\newcommand{\partialto}{\to}

\newcommand{\BBtheFifth}{47{,}176{,}870}
\newcommand{\BBtheFifthTNF}{181{,}385{,}789}
\newcommand{\SigmaTheFifth}{4{,}098}

\newcommand{\BBTxF}{3{,}932{,}964}
\newcommand{\BBTxFTNF}{2{,}154{,}217}

\newcommand{\numBBsholdouts}{1{,}214}

\newcommand{\radofull}{Tibor Rad\'o\xspace}
\newcommand{\rado}{Rad\'o\xspace}

\newcommand{\cryptid}{Cryptid\xspace}

\newcommand{\cycler}{Cycler\xspace}
\newcommand{\cyclers}{Cyclers\xspace}
\newcommand{\TC}{Translated Cycler\xspace}
\newcommand{\TCs}{Translated Cyclers\xspace}

\newcommand{\headpos}{head-position\xspace}
\newcommand{\headposs}{head-positions\xspace}

\newcommand{\states}{\mathcal{S}}
\newcommand{\alphabet}{\mathcal{A}}
\newcommand{\balphabet}{\left\{\szero,\sone\right\}}
\newcommand{\symbolzero}{\texttt{0}}
\newcommand{\symbolone}{\texttt{1}}

\newcommand{\numSporadic}{13\xspace}

\newcommand{\ssp}{state-symbol pair\xspace}
\newcommand{\ssps}{state-symbol pairs\xspace}

\newcommand{\HALT}{\texttt{HALT}\xspace}
\newcommand{\NONHALT}{\texttt{NONHALT}\xspace}
\newcommand{\UNKNOWN}{\texttt{UNKNOWN}\xspace}

\newcommand{\CoqBB}{Coq-BB5\xspace}
\newcommand{\CoqBBnospace}{Coq-BB5}

\newcommand{\TMstep}{\to}

\usepackage{thm-restate}
\usepackage{makecell}

\newcommand{\topleftsymbolcell}[3]{% color, symbol, content
    \cellcolor{#1}%
    \begin{tikzpicture}[overlay, remember picture]
        \node[anchor=north west, inner sep=1pt] at (0,1.7ex) {\small #2};
    \end{tikzpicture}%
    \hspace{1.5ex}#3%
}

\begin{document}

%%
%% Title
%%
\title{Determination of the fifth Busy Beaver value}

%%
%% Authors - TODO: Add affiliations and emails for each author.
%% Each author gets their own \author{} block in acmart.
%%
%% TODO: Fill in correct affiliations, emails, and countries for all authors.
% \author{The bbchallenge Collaboration}
% \affiliation{\institution{bbchallenge.org}\country{}}

\author{Justin Blanchard}
\authornote{Alphabetical ordering of authors after the collaboration name.}
\affiliation{\institution{}\country{USA}}

\author{Daniel Briggs}
\affiliation{\institution{}\country{USA}}

\author{Konrad Deka}
\affiliation{\institution{}\country{Poland}}

\author{Nathan Fenner}
\affiliation{\institution{}\country{USA}}

\author{Yannick Forster}
\affiliation{\institution{INRIA Paris}\country{France}}

\author{Georgi Georgiev (Skelet)}
\affiliation{\institution{Sofia University, Faculty of Mathematics and Informatics}\country{Bulgaria}}

\author{Matthew L. House}
\affiliation{\institution{}\country{USA}}

% \author{Rachel Hunter}
% \affiliation{\institution{}\country{}}

% \author{Iijil}
% \affiliation{\institution{}\country{}}

\author{Maja K\k{a}dzio{\l}ka}
\affiliation{\institution{University of Warsaw}\country{Poland}}

\author{Pavel Kropitz}
\affiliation{\institution{}\country{Slovakia}}

\author{Shawn Ligocki}
\affiliation{\institution{}\country{USA}}

% \author{mxdys}
% \affiliation{\institution{}\country{}}

\author{Mateusz Na\'{s}ciszewski}
\affiliation{\institution{}\country{Poland}}

% \author{savask}
% \affiliation{\institution{}\country{}}

\author{Tristan St\'erin}
\authornote{Corresponding author.}
\affiliation{\institution{PRGM DEV}\country{France}}
\email{tristan@prgm.dev}

\author{Chris Xu}
\affiliation{\institution{UC San Diego}\country{USA}}

\author{Jason Yuen}
\affiliation{\institution{}\country{Canada}}

\author{Th\'eo Zimmermann}
\affiliation{\institution{LTCI, T\'el\'ecom Paris, Institut Polytechnique de Paris}\country{France}}

%%
%% Abstract
%%
\begin{abstract}
    The Busy Beaver value $S(n)$ is the maximum~number of steps that an n-state 2-symbol Turing machine can perform from the all-zero tape before halting. $S$ was historically introduced by \radofull in 1962 as one of the simplest examples of an uncomputable function.

    We prove that $S(5) = 47,176,870$ using the Coq proof assistant, confirming the conjecture of Aaronson \cite{BusyBeaverFrontier}. The proof enumerates 181,385,789 Turing machines with 5 states and, for each machine, decides whether it halts or not.
    Our result marks the first determination of a new Busy Beaver value in over 40 years and the first Busy Beaver value ever to be formally verified, attesting to the effectiveness of massively collaborative online research (\url{bbchallenge.org}).
\end{abstract}

\maketitle


Blablabla





\bibliographystyle{ACM-Reference-Format}
%\newpage
\bibliography{bbchallenge-paper}







\end{document}
